\begin{lemma}[Step 1, Subcase 1.1: the later end is outside the starred predecessor]
Let $G$ be a finite loopless multigraph with individual edge identities,
$e_0$ a distinguished edge, and $[k]$ the color set.  Let $T$ be a
positive-rung ETT with generating coloring $\varphi$, maximum property,
and a certified Tashkinov series.  Assume that $\varphi$ is proper and
leaves exactly $e_0$ uncolored.  Assume the complete lower-rung package:
for every positive-rung maximum ETT of smaller rung, its terminal vertex
set is elementary, its final ladder closure has the stated stable
interchangeability property, PE stability descends and equal-level PE
stored pairs are disjoint, a terminal PE component meets its ladder tree
only at its support, nonminimum defective endpoints precede the later
allowed PE/SE vertices, and every terminal stable coloring is a
mod-coloring.

Let $B$ be a hierarchy base for $T$, let $H$ be a good hierarchy with
$d$ dividers, and fix an integer $q_0>d$.  Assume uniformly that every
same-rung maximum ETT, built from an arbitrary certified input and
generating coloring and carrying a good hierarchy with fewer than $q_0$
dividers, is elementary.  Assume also the large-branch estimate
$|\overline\varphi(T_n)|\geq 2n+11$.  Write
$A=H.\mathit{starred}(d)$ for the final starred predecessor.  Besides all
five fields of the good hierarchy, assume separately that the terminal
tree is
\[
 \left(\bigcup_{\eta\in D^H_{n,d+1}}\Gamma^d_\eta\right)^-
 \text{-closed under }\varphi,
\]
and assume that $V(T)$ is elementary under $\varphi$.

Let $\alpha\ne\beta$, with $\alpha$ missing and singleton-closed on $T$
and $\beta$ missing on $T$.  Let $P_1,P_2$ be exact
$(T,\varphi,\{\alpha,\beta\})$-exit paths, with respective outside ends
$a_1,a_2$, tree ends $b_1,b_2$, disjoint incidence-walk supports, and
distinct outside and tree ends.  Order them so that $b_1\prec_T b_2$.
For each $j$, retain the exact edge identity $f_j$ on $P_j$ incident
with $b_j$, and assume that $f_j$ is a $\beta$-colored boundary edge of
$T$.  If $b_2\notin V(A)$, then these data are contradictory.
\end{lemma}
\begin{proof}
Put
$C_\Gamma=\bigcup_{\eta\in D^H_{n,d+1}}\Gamma^d_\eta$.
The exact boundary identity $f_2$ has color $\beta$.  If $\beta$ were
outside $C_\Gamma$, the stated terminal
\noderef{TreeColorSetMinusClosed} certificate, together with the fact
that $\beta$ is missing on $T$, would make its boundary empty.  Hence
$\beta\in\Gamma^d_\eta$ for some
$\eta\in D^H_{n,d+1}$.  Unfolding
\noderef{TashkinovStoredColors},
\noderef{HierarchyStoredColorsAt}, and
\noderef{HierarchyTerminalStoredColors}, and using $A\subseteq T$, gives
$D^H_{n,d+1}\subseteq D^H_{n,d}$.  The exact preceding-stage conclusion
of \noderef{GoodHierarchyReservedDisjoint} therefore makes $\beta$
missing on $A$.

Choose a color $\gamma$ missing at $b_2$.  The choice is not an extra
assumption: the initial constructor of \noderef{TashkinovSeriesStep}
contains a \noderef{KTriple}, so $k\geq\Delta(G)+1$, whereas properness,
\noderef{Degree}, and \noderef{MaximumDegree} give at most
$d_G(b_2)\leq\Delta(G)$ present colors.  Thus the complement in
\noderef{MissingColors} is nonempty.  If $\gamma$ belonged to
$C_\Gamma$, \noderef{GoodHierarchyReservedDisjoint} would make it
missing at a vertex of $A$.  That vertex differs from $b_2$ because
$b_2\notin V(A)$, contradicting the assumed
\noderef{ColoredVertexSetElementary} property of $V(T)$.  Hence
$\gamma\notin C_\Gamma$, and the terminal minus-closed certificate makes
$\gamma$ closed on $T$.  The color $\alpha$ is closed there by
hypothesis.  Since $f_2$ is a $\beta$-edge incident with $b_2$, while
$\gamma$ is missing at $b_2$, we also have $\gamma\ne\beta$.

If $\alpha=\gamma$, set $\mu_1=\varphi$ and retain the original certified
ETT, maximum property, hierarchy, and terminal certificate.  Otherwise
exchange $\alpha$ and $\gamma$ on every full two-color component outside
$T$.  The two singleton closedness assertions show that this is an
\noderef{OutsideKempeRecoloring}.  By
\noderef{OutsideKempeRecoloringStable}, $\mu_1$ is proper, leaves exactly
$e_0$ uncolored, fixes every individual edge incident with $T$, and is
stable on every displayed stage.  Applying \noderef{HierarchyLEMStable}
with the stated lower package gives a corresponding same-rung ETT under
$\mu_1$, maximum property, the identical ordered hierarchy and Gamma
family, all five good-hierarchy fields, and the transported terminal
minus-closed certificate.  This lists the replacement data explicitly;
none is inferred merely from stability.

Because $\alpha$ and $\gamma$ are both closed on $T$, every external
$\alpha$-edge of $P_1$ or $P_2$ lies in one of the exchanged components.
It becomes a $\gamma$-edge.  Every $\beta$-edge is fixed, using both
$\beta\ne\alpha$ and $\beta\ne\gamma$.  In the equality branch the paths
are unchanged.  Thus the same exact edge-node walks $P_1,P_2$ are
$(\gamma,\beta)$-exit paths under $\mu_1$.

Let $K_2$ be the full $(\gamma,\beta)$ component through $b_2$, and let
$\mu_2$ be obtained by swapping that component.  The vertex $b_2$ misses
$\gamma$ and is incident with the $\beta$-edge $f_2$, so it has incidence
degree one.  The outside end $a_2$ misses one path color by
\noderef{TreeExitPath} and is incident with the other along the nontrivial
path, so it also has degree one.  By
\noderef{TwoColorIncidenceDegreeAtMostTwo} and
\noderef{TwoColorComponentSpanningPath}, no further node can attach to
$P_2$: attachment internally would create degree three, while attachment
at an endpoint would raise its degree from one.  Therefore $P_2$ is the
whole component $K_2$.  Its exit-path condition says that it meets $T$
only at $b_2$; in particular it is disjoint from the exact ordered prefix
$T^-$ immediately before $b_2$.

The colors of the swapped component are distinct.  Hence
\noderef{TwoColorComponentSwapProper} makes $\mu_2$ proper, and
\noderef{TwoColorComponentSwap} changes only two nonempty colors, so the
unique uncolored identity remains exactly $e_0$.  Repackage the hierarchy
on $T^-$ as follows.  If $A\subsetneq T^-$, retain all $d$ dividers and
use $T^-$ as terminal.  If $T^-=A$ and $d>0$, first use the valid
$(d-1)$-divider presentation ending at $A$.  If $T^-=A$ and $d=0$, use
the actual initial-tuple route through \noderef{TechnicalTheoremBase}; no
zero-rung ETT is constructed.  The exact finite closure sequence in
\noderef{TashkinovClosure} and
\noderef{TashkinovAugmentationSequence} restricts to the retained prefix.
The mod-coloring conclusion of \noderef{Td} then permits the single exact
edge-and-vertex extension through $b_2$.  Thus
\noderef{CorrespondingExtendedTashkinovTree} and
\noderef{MaximumPropertyTransport} produce the exact prefix $C$ through
$b_2$ as a corresponding same-rung maximum ETT.

On $C$, copy every old stage through $A$.  The last containment is now
strict because the individual edge and vertex at $b_2$ were restored.
The stored sets, pair sizes, disjointness, unique refresh, parallel-base
field, and internal minus-closed fields are unchanged; terminal avoidance
is checked on that exact restored edge.  Hence this is a genuine good
hierarchy with $d<q_0$ dividers, not a record asserted by name.

The component swap is disjoint from $T^-$, so the earlier missing witness
for $\beta$ in $A\subseteq T^-$ remains a missing witness under $\mu_2$.
The swap makes $\beta$ missing at $b_2$.  These are two distinct vertices
of $C$.  The same-rung fewer-divider hypothesis applies to the newly
constructed arbitrary-input, corresponding maximum ETT and its verified
good hierarchy, and makes $V(C)$ elementary under $\mu_2$.  The two
$\beta$-missing vertices contradict this elementarity.  Therefore the
data of Subcase 1.1 cannot occur.
\end{proof}
